\chapter{模板template}
Django 模板是一个文本文件，它包含静态的 HTML，以及嵌入式的 Django 模板语言(DTL)代码。

\section{模板概述}
在Django中，创建一个\codeinline{base.html}文件，作为模板的基础文件，其他模板文件可以继承这个基础文件，从而减少重复代码。

\begin{enumerate}[label=\arabic*)]
  \item 一般在\textcolor{YBXPurple}{主项目根目录}创建一个\codeinline{templates}文件夹，用于存放模板文件。
\begin{lstlisting}[style=demostyle,numbers=none]
$> tree 
.
|-- manage.py
|-- TestModel  # 子应用
|   |-- admin.py
|   |-- apps.py
|   |-- views.py
|   |-- urls.py
|   |-- models.py
|   |-- tests.py
|   `-- migrate
|-- mysite   # 主项目目录
│   |-- __init__.py
│   |-- settings.py
│   |-- urls.py
│   `-- wsgi.py
`-- templates # 模板文件夹
\end{lstlisting}

  \item 在\textcolor{YBXPurple}{主项目}的\codeinline{settings.py}文件中配置模板文件路径。
    
\begin{lstlisting}[style=demostyle,numbers=none]
$> cat mysite/settings.py    
TEMPLATES = [
    {
        'BACKEND': 'django.template.backends.django.DjangoTemplates',
        # 'DIRS': [], # 之前的配置
        # 修改为
        'DIRS': [os.path.join(BASE_DIR, 'templates')],       
        'APP_DIRS': True,
        'OPTIONS': {
            'context_processors': [
                'django.template.context_processors.debug',
                'django.template.context_processors.request',
                'django.contrib.auth.context_processors.auth',
                'django.contrib.messages.context_processors.messages',
            ],
        },
    },
]
\end{lstlisting}

  \item 在\codeinline{templates}文件夹下创建一个\codeinline{base.html}文件，作为模板的基础文件，其他模板文件可以继承这个基础文件，从而减少重复代码。
\begin{lstlisting}[style=demostyle, numbers=none]
$> cat templates/base.html
<!DOCTYPE html>
<html lang="en">
<head>
    {% block delihead %}
    <meta charset="UTF-8">
    <meta name="viewport" content="width=device-width, initial-scale=1.0">
    <title>测试文件标题</title>
    <link rel="icon" href="static/images/logo.png">
    <link rel="stylesheet" href="static/style.css" type="text/css">
    {% endblock %}
</head>

<body>
    <h2>美使馆</h2>
    {% block content %}
    {% endblock %}
</body> 

<footer>
    <small>&copy 2020 美使馆<a href="http://www.baidu.com">JJW</a></small>
</footer>

</html>
\end{lstlisting}

  \item 在\codeinline{urls.py}文件中配置路径。
\begin{lstlisting}[style=demostyle,numbers=none]
$> cat urls.py
from django.urls import path
 
from . import views
 
urlpatterns = [
    path('hello/', views.hello),
]
\end{lstlisting}

\item 在\textcolor{YBXPurple}{子应用}的目录可以创建一个\codeinline{templates}文件夹，用于存放子应用的模板文件。
\begin{lstlisting}[style=demostyle,numbers=none]
$> tree TestModel
TestModel/
|-- __init__.py
|-- admin.py
|-- apps.py
|-- __pycache__
│   |-- __init__.cpython-39.pyc
│   |-- views.cpython-39.pyc
│   `-- urls.cpython-39.pyc
|-- templates
│   `-- hello.html
|--tests.py
|-- urls.py
`-- views.py
\end{lstlisting}  

其中 \codeinline{hello.html} 是子应用\codeinline{TestModel}的模板文件。
\begin{lstlisting}[style=demostyle,numbers=none]
$> cat templates/hello.html
{% extends 'base.html' %} <!-- 继承temples的base.html -->

{% block content %} <!-- 重写base.html的content -->
    <h3>Register</h3>
    <form action="/login/create_user" method="POST">
        {% csrf_token %}
        <input type="text" name="username" placeholder="Username" required>
        <input type="email" name="email" placeholder="Email" required>
        <input type="password" name="password" placeholder="Password" required>
        <input type="password" name="password2" placeholder="Confirm Password" required>
        <button type="submit">Register</button>
    </form>
{% endblock %}
\end{lstlisting}
  
  \item 在\textcolor{YBXPurple}{子应用}的\codeinline{views.py}文件中增加一个新的对象，
  用于向模板提交数据。
\begin{lstlisting}[style=demostyle,numbers=none]
$> cat TestModel/views.py
from django.shortcuts import render

"""
一个视图函数 hello
  接受一个 request 对象作为参数，并返回一个 HTTP 响应
"""
def hello(request):  
    context = {}     
    # context 是一个字典，用来存储传递给模板的上下文数据。在这个字典中，键是模板变量的名称，值是要传递给模板的实际数据。
    context['hello'] = 'Hello World!' 
    # 将 'Hello World!' 赋值给 context 字典中的 'hello' 键。这样，hello 将作为模板变量传递给模板，并且在模板中可以使用这个值。
    return render(request, 'hello.html', context)
\end{lstlisting}

  \item 启动服务器，访问\codeinline{http://localhost:8000/hello/}，可以看到页面显示\codeinline{Hello World!}。

\end{enumerate}

\subsection{模板的定义}
\begin{enumerate}[label=\arabic*)]
    \item 模板变量：\codeinline{\{\{ 变量名 \}\}}
    \item 从URL获取模板变量
\begin{lstlisting}[style=demostyle, numbers=none]
{% url '路由名' 参数 %}    
\end{lstlisting}    

\begin{example}
    \exProblem
    在\codeinline{urls.py}中定义一个
\begin{lstlisting}[style=demostyle, numbers=none]
re_path('(?P<year>[0-9]{4}).html', views.myyear, name='urlyear'),        

# 1.(?P<year>[0-9]{4})：
这是正则表达式中的一个命名捕获组，(?P<year>[0-9]{4})的含义是匹配4位数字（0-9之间），并将匹配到的内容命名为year。
# 2.views.myyear：当URL匹配成功后，Django会调用 views.myyear 这个视图函数，并将匹配到的 year 参数传递给它。
例如，如果URL是 2023.html，那么 views.myyear 会接收到 year=2023。
# 3.name='urlyear'：
这是给这个URL模式起的一个名字，方便在Django的其他地方通过名称来引用该URL。
通过这个名字，你可以使用 reverse() 函数或 {% url %} 模板标签来生成对应的URL，例如 {% url 'urlyear' year=2023 %} 会生成 2023.html。
\end{lstlisting}    


    
而在模板文件夹的\codeinline{test.html}中，
\begin{lstlisting}[style=demostyle, numbers=none]
<a href="{% url 'urlyear' 2019 %}">2019年度</a>
\end{lstlisting}        
    \exSolution

    当找到\codeinline{test.html}模板文件时，会到\codeinline{urls.py}中找到\codeinline{urlyear}这个路由，
    并将\codeinline{2019}作为参数传递给\codeinline{myyear}函数。
    从而在浏览器中显示\codeinline{2019年度}的超链接。

\begin{lstlisting}[style=demostyle, numbers=none]
<a href="/2019.html">2019年度</a>
\end{lstlisting}            
    
\end{example}

    \item 读取静态资源内容
\begin{lstlisting}[style=demostyle, numbers=none]
{% static "css/style.css" %} 
\end{lstlisting}
    \item for遍历：\codeinline{\{\% for 变量名 in 列表名 \%\}}
\begin{lstlisting}[style=demostyle, numbers=none]
{% for item in list %}
    {{ item }}
{% endfor %}    
\end{lstlisting}
    \item if判断：\codeinline{\{\% if 条件 \%\}}
\begin{lstlisting}[style=demostyle, numbers=none]
{% if 条件 %}
    {{ 变量 }}
{% endif %}
\end{lstlisting}       
\end{enumerate}


\section{静态文件}

\subsection{静态文件的作用}

静态文件是指不需要经过服务器处理，直接可以被浏览器访问的文件，例如：图片、CSS、JavaScript等。

\begin{enumerate}[label=\arabic*)]
    \item 图片：\codeinline{.jpg}、\codeinline{.png}、\codeinline{.gif}等。
    \item CSS：\codeinline{.css}文件，用于设置网页的样式。
    \item JavaScript：\codeinline{.js}文件，用于设置网页的交互效果。
\end{enumerate}

\subsection{静态文件的配置}
\begin{enumerate}[label=\arabic*)]
  \item 在 \textcolor{YBXPurple}{主项目}的\codeinline{settings.py}文件中配置静态文件路径。
\begin{lstlisting}[style=demostyle,numbers=none]
STATIC_URL = '/static/'
if not DEBUG:
    STATIC_ROOT = os.path.join(BASE_DIR, 'static')

# 重写 STATICFILES_DIRS    
STATICFILES_DIRS = (os.path.join(BASE_DIR, 'static'),) # 必须是一个元组(,)要有,
\end{lstlisting} 

  \item 想要引用static资源文件的\textcolor{YBXPurple}{子项目}，
  需要在\textcolor{YBXPurple}{子项目}的\codeinline{urls.py}文件中的\codeinline{urlpatterns}列表中添加一条配置，
  用于访问静态文件。

\begin{lstlisting}[style=demostyle,numbers=none]
from django.urls import path
from . import views  # Import the views module

from django.conf import settings
from django.conf.urls.static import static

urlpatterns = [
    path("", views.open_register, name="open register"),  
    
] + static(settings.STATIC_URL, document_root=settings.STATICFILES_DIRS[0]) 
# 在 settings.py 中配置了 STATICFILES_DIRS，所以这里可以直接使用 settings.STATICFILES_DIRS[0] 来获取静态文件目录
\end{lstlisting} 
\end{enumerate}

在\textcolor{YBXPurple}{子项目}的视图\codeinline{views.py}中，可以通过\codeinline{render}函数将静态文件传递给模板。
\begin{lstlisting}[style=demostyle,numbers=none]
from django.shortcuts import render
from django.template import loader
from django.http import HttpResponse

# Create your views here.
def open_register(request):
    template = loader.get_template('login/register.html')
    return HttpResponse(template.render({}, request))
\end{lstlisting}






